% Cote 119: offprints, reports, courses, CV and letters (typed and printed part): archive cote 119, site page 222 (PDF page 223).
% typescript. Draft compared: none.
% AI-transcribed by Claude Opus 5.5 (claude-opus-5-5, adaptive thinking, effort high), Message Batches API, batch r119_full (msgbatch_01DwZNm5Y7YsKBWyAN4sFs8w), 2026-11-13; not yet checked by eye.
\scanpage{119/222}
\mspage{. 9 .}

6.\quad \uline{Groupes algébriques} ([16 SGA 3 - en trois volumes]
[12 bis]).

Ce sujet relève à la fois de la géométrie
algébrique et de la théorie des groupes. Le travail
cité SGA 3 se place surtout sur des schémas de base
généraux, et la part de la géométrie algébrique y est
certes considérablement plus large que celle de la
théorie des groupes. Néanmoins, grâce à la technique
des schémas, nous y obtenons des résultats nouveaux
même dans le cas de groupes définis sur un \uline{corps} de
base, les plus intéressants (relatifs surtout au cas
d'un corps de base imparfait) étant contenus dans
SGA 3, Exp XIV. Ma contribution principale, continuant
dans la voie ouverte par A. Borel et C. Chevalley dans
le contexte de la géométrie algébrique habituelle, a
été de montrer le parti qu'on pouvait tirer d'une
application systématique de la théorie des schémas aux
groupes algébriques et aux schémas en groupes.

\uline{Références} : J. Tits, F. Bruhat, M. Demazure, P. Gabriel,
A. Borel, D. Mumford.


7.\quad \uline{Groupes discrets} ([18, Exp VIII], [13 bis]).

Dans [18, Exp. VIII] je développe une théorie
purement algébrique des classes de Chern des représen-%
tations d'un groupe discret sur un corps de base
(ou même un anneau de base) quelconque, avec des appli-%
cations de nature arithmétique sur l'ordre des classes
de Chern des représentations complexes. Cette théorie
peut être considérée comme cas particulier d'une théorie
des classes de Chern des représentations linéaires
de schémas en groupes quelconques, elle-même contenue
dans la théorie des classes de Chern $\ell$-adiques des fibrés

\margin{\unsure{222|1222}}
